\begin{abstract}
Edge object detectors deployed on passively cooled single-board computers
fail in ways a fixed model cannot absorb: sustained load drives the SoC into
firmware throttling, frame rate becomes erratic, and a crashed or hung process
goes unnoticed. We present \shade, an open-source control plane that keeps a
detector within a thermal and memory envelope by moving along a
\emph{degradation ladder} of operating points (model, input size, detection
interval), with hysteresis-guarded one-step recovery, background model
hot-swap, a supervising watchdog, and an audit log that records why every
action was taken. On top of the ladder we add a content-aware scheduler that
spends the ladder's detection budget where the scene changes: it stretches the
detection interval while the scene is static and moves previous boxes with the
scene by phase correlation on skipped frames. In an offline study with a real
YOLOv8n model, gating removed \simGateSaving\% of detector calls on a clip that is
half static with no change in agreement, and box shifting raised F1 at IoU~0.75
under fast camera motion from \simShiftGain{} (synthetic pan of one image).
On a Raspberry Pi~5 (\TBD{board, cooling, enclosure}) \shade{} reduced
firmware-throttled time from \resStaticThrottled\% to \resShadeThrottled\% and raised
the 5th-percentile frame rate from \resStaticPfive{} to \resShadePfive{}~FPS
relative to a static nominal configuration, at a mean-frame-rate cost of
\TBD{X\%}. We also report negative results from a closed-loop simulator:
a learned thermal-forecast policy did not outperform tuned reactive
control, and allowing the firmware to throttle yields higher mean throughput
than any ladder unless throttling is costly, which delimits when adaptation
pays. Code, data and scripts are released.
\end{abstract}
